\documentclass[11pt]{article}
\usepackage[margin=1in]{geometry}
\usepackage{amsmath,amssymb}
\usepackage[hidelinks]{hyperref}

\begin{document}

\section*{Human Review: Borel Factor-Subspace Selection}

\noindent Original automatic proof: \texttt{solution\_packet.pdf}

\noindent Checked date: 19 June 2026

\noindent Status: Partially manually checked; author verification recommended.

\noindent Verdict: The proof looks right, but I am not able to verify every technical detail to the appropriate level.

\section*{Human Comments}

The proposed argument addresses Question 6.1 of Antunes--Beanland--Braga,
\emph{Uniformly factoring weakly compact operators and parametrised
dualisation}.  The question asks whether the Borel measurable versions of
Theorem 5.1, Corollary 5.12, and Corollary 5.13 hold.

The proof appears to make one precise repair to the published proof of
Theorem 5.1.  In the original argument, after constructing the Borel code map
\[
  \eta=\psi\circ\sigma:\mathcal B\longrightarrow \mathbb N^{\mathbb N},
\]
the authors choose a closed tree whose projection is the analytic set
\(\eta(\mathcal B)\).  A branch above each code is then chosen using
Jankov--von Neumann uniformization, which gives only
\(\sigma(\Sigma^1_1)\)-measurability.  The proposed proof observes that
\(\eta(\mathcal B)\) is not an arbitrary analytic set: it is the image of a
Borel map.  It therefore builds the closed tree from a closed parametrization
of the Borel graph of \(\eta\), obtaining a genuinely Borel branch selector.

The central descriptive-set-theoretic lemma looks correct.  Given a Borel map
\(\eta:S\to\mathbb N^{\mathbb N}\), one first codes the standard Borel space
\(S\) as a Borel subset of Baire space, considers the Borel graph
\[
  G=\{(i(s),\eta(s)):s\in S\},
\]
and represents this Borel set as a one-to-one continuous image of a closed
set.  This gives a closed set
\[
  C\subseteq\mathbb N^{\mathbb N}\times\mathbb N^{\mathbb N}
\]
and a Borel lift \(\tau:S\to C\) with \(p(\tau(s))=\eta(s)\).  This seems to
supply exactly the missing Borel branch choice.

The key structural check also appears to work.  Every branch of the new tree
still has first coordinate equal to \(\eta(A_0)=\psi(\sigma(A_0))\) for some
original parameter \(A_0\in\mathcal B\).  Thus the modified tree should not
introduce any new ``bad'' branch spaces.  The reflexivity argument in the
Kurka amalgamation step is therefore the same as in the source proof: the
branch spaces are reflexive, and Kurka's proposition then gives the
reflexivity of the amalgamated space \(Z\).

The measurability repair also appears sound.  In the source proof, the proof
of Claim 5.11 explicitly reduces the measurability of \(\Phi\) to the
measurability of the assignment \(A\mapsto t_A\), together with already Borel
objects such as \(\sigma(A)\), \(\varphi_A\), and the fixed selectors.  Replacing
the old \(\sigma(\Sigma^1_1)\)-measurable selector by the Borel selector
\(t_A=\tau(A)\) should therefore upgrade the same displayed formula from
\(\sigma(\Sigma^1_1)\)-measurability to Borel measurability.  The Borelness of
\(\Psi(A)=\overline{\operatorname{span}}\{\Phi(A,d_n(X)):n\in\mathbb N\}\)
then follows from the usual Kuratowski--Ryll-Nardzewski selectors and the
standard Borel closed-span operation.

\section*{Citation Correction}

There is a small citation-numbering correction in the automatic proof.  The
rational-code map
\[
  \psi:\mathrm{mbs}\to\mathbb N^{\mathbb N}
\]
should be cited as coming from Lemma 5.9 of Antunes--Beanland--Braga, not
Lemma 5.10.  In the source paper, Lemma 5.9 is the result defining and
recording the properties of \(\psi\); Claim/Lemma 5.10 in that proof is the
later reflexivity claim for the interpolation space \(Z\).

\section*{Caveat}

Although the method appears solid and the single proposed modification seems
to target exactly the obstruction in the published proof, I am not sufficiently
familiar with all the descriptive-set-theoretic and parametrized Banach-space
machinery involved to certify the proof at the level one would want before
making a firm mathematical claim.  In particular, small technical issues could
still be hidden in the passage from the modified closed tree to the exact
Kurka amalgamation framework, or in the precise Borel coding of the resulting
maps.

\section*{Review Outcome}

This should be treated as a promising and apparently correct solution, but not
as fully human-certified at this stage.  I recommend sending the argument to
the authors to ask whether the Borel graph/closed-tree
replacement indeed removes the \(\sigma(\Sigma^1_1)\)-measurability obstacle in
Question 6.1.

\end{document}

